\chapter{Future work}
\label{ch:future}

Ordered by what would change a conclusion in this report, rather than by what
would be pleasant to build.

\section{Measure something real}

\textbf{The single most valuable next measurement is a small, hand labelled set
of real pages.} Every accuracy number in this report is a measurement on a
synthetic corpus, and the strongest objection to all of them is that a generated
corpus has a generator's habits.

This does not require scraping, which is out of scope and stays out of scope.
It requires a small set of pages, labelled by hand against the token set, with
the labelling protocol and the disagreement rate published alongside. A hundred
real forms with a published inter annotator agreement would tell a reader more
about whether the tool works than another thousand generated ones.

The prediction that goes with it should be registered first, and the honest
registration is that the rule table's advantage will shrink, because a
substantial part of it is that its vocabulary and the corpus's label text were
written by the same person.

\section{Fix the double classification in the evaluation runner}

Described in Chapter~\ref{ch:limitations}. The runner classifies each form twice
and the two passes are separate draws for a non deterministic engine, which puts
an excess disagreement on the order of two fields in seven hundred between a
row's predicted label and its finding codes.

The fix is to pass the first pass's predictions into the audit rather than
letting the audit reclassify. It is small, it makes the two numbers one fact, and
it halves the cost of any future language model benchmark. It was not done in the
phase that found it because changing what the runner does for every engine after
seeing the results is exactly the regeneration this project's rules forbid, and
because the honest place for it is a phase that re-registers its predictions and
re-runs everything.

\section{Vary the one thing that was held fixed}

The language model result rests on one model, one quantization and one prompt.
Three variations would each answer a different question, and each needs its own
pre-registration:

\begin{itemize}
  \item \textbf{A higher precision quantization of the same model.} Isolates the
    quantization from the model. Cheap to run and answers an objection that will
    otherwise be raised forever.
  \item \textbf{A second prompt, registered before it is run.} A prompt sweep
    chosen after seeing a result is a search for a number; a second prompt
    registered in advance, with a stated hypothesis about which slice it should
    move, is a measurement.
  \item \textbf{A fine tuned small model.} The comparison as it stands is
    between two engines fitted on the corpus and one that met it at inference.
    Removing that asymmetry is the fairest version of the experiment and it is
    also the most expensive.
\end{itemize}

\section{The transformer rung}

The classifier ladder specifies a small multilingual encoder, fine tuned and
exported with dynamic integer quantization, as an optional fourth engine
\cite{onnxruntime-quantization}. It has not been built, which is why the fourth
row of Table~\ref{tab:headline} reads \emph{absent} in words.

The condition for shipping it is stated in advance and is not negotiable
afterwards: it ships only if it beats a pre-stated metric by a pre-stated margin
within a pre-stated latency budget. If it does not, the deliverable is the
paragraph saying so, with the numbers, and that is a successful outcome rather
than a failed one.

Two things the tier profile in Figure~\ref{fig:tier-profile} suggests it should
be measured against specifically. The mixed tier, where the n-gram model
underperforms its own tier by tier behaviour and where a model with more context
might do better. And the held out locale, where a multilingual encoder's
pretraining is the whole argument for using one.

\section{More templates per family}

The design resolves large effects and not small ones, and the only thing that
improves that is more templates in the test partition. The arithmetic is in
Table~\ref{tab:power-repair}: the relationship between templates and the smallest
attainable p value is steep at the low end and flattens quickly.

The cost is the one already paid once: a regeneration, a retraining, a threshold
rederivation and a re-measurement, with every recorded number moving at once and
an explanation owed in advance.

\section{Close the cross repository tasks}

Five issues are open on the evaluation harness, each with a concrete API
proposal, and two of them block things this project would like to state
plainly:

\begin{itemize}
  \item \textbf{Publication of the harness to a package index.} Until then the
    dependency is a source reference rather than a released version, and the
    project's own rule forbids tagging a stable release while that is true. This
    is the one open task that is a hard blocker rather than an improvement.
  \item \textbf{A narrower dependency set for the harness.} It currently pulls
    more than a dozen transitive packages including a training metrics logging
    stack, which is why it sits in the development extra and why the shipped tool
    cannot compute its own significance tests.
  \item Clustered and paired resampling as a public entry point, an absolute
    practical effect threshold, a stable join key on the verdict output, and a
    typing marker so that values crossing the boundary are checked.
\end{itemize}

\section{Smaller things worth doing}

\begin{itemize}
  \item \textbf{Make the Japanese postal mark reachable.} It is in the rule table
    with a comment saying it cannot fire, because normalisation treats a lone
    symbol as a delimiter. Making it reachable means changing what counts as a
    token character, which churns every golden snapshot and every committed
    fixture expectation, so it is a deliberate change with a cost rather than a
    quiet fix.
  \item \textbf{Ship a model with the wheel, or a documented way to fetch one.}
    A published wheel currently carries no model, so an install runs the rule
    baseline until the user trains one or points the tool at a bundle. The
    behaviour is documented and prints a line rather than failing silently, but
    it means the default install is not the default engine the measurements
    recommend.
  \item \textbf{Widen the hyperparameter grid, with a fresh registration.} The
    sweep chose the edge of its own grid, so the shipped model may be less
    regularised than a wider grid would have chosen. Widening it now would be
    searching for a number; widening it inside a phase that registers first is a
    measurement.
  \item \textbf{Report the reliability curve per class rather than pooled.}
    Calibration made the pooled expected calibration error worse while improving
    macro-F1, and a per class view would say which classes moved and in which
    direction.
\end{itemize}

\section{What is deliberately not future work}

Named so that they stay named and stay out.

Crawling. Form filling. Markup rewriting. Scraping training data. A browser
extension that duplicates the extractor in a second language and doubles the
surface that has to be kept in agreement. Each of those is a boundary rather than
a missing feature, and each is stated as a boundary in the project's own
documentation.

A form filling agent is a genuinely interesting successor and it is named here so
that it stays a separate project rather than becoming a flag on this one.
